\section{Real remainders from a simultaneous Pad\'e construction}\label{sec:pade}
The preceding construction begins with many beta values and tries to
remove some of them. The present family has a remainder involving
$\beta(4)$ alone from the start. It comes from the $m=5$ construction
of Kawashima--Po\"els~\cite{KP}, which approximates two formal functions
with a common denominator polynomial. The two remainders are separate;
simultaneous approximation does not introduce a $\beta(2)$ term into
the weight-four remainder.

The first task is to give that formal remainder an actual real meaning.
The second is to estimate it when both the degree and evaluation point
grow. Formal order of approximation at a fixed degree does not answer
the latter question.

\subsection{Functions and polynomials}
For $s=2,4$ and $\Re z>0$, set
\[
 F_s(z)=(s-1)!\sum_{j\ge0}\frac{(-1)^j}{(z+j)^s}
       =\int_0^\infty\frac{u^{s-1}e^{-zu}}{1+e^{-u}}\dd u.
\]
The shift identity $F_s(z)+F_s(z+1)=(s-1)!z^{-s}$ determines its
formal asymptotic Laurent coefficients $c_{s,j}$ by
\begin{equation}\label{eq:pade-moments}
 2c_{s,j}=(s-1)!\,[j=s-1]
              -\sum_{v<j}(-1)^{j-v}\binom jv c_{s,v}.
\end{equation}
The indexing convention is $F_s(z)\sim\sum_{j\ge0}c_{s,j}z^{-j-1}$;
call the corresponding formal series $\Theta_s$. In
\eqref{eq:pade-moments}, $[j=s-1]$ is one when $j=s-1$ and zero otherwise.
A finite Taylor expansion of the integrand near zero, with a separate
exponentially small tail, proves that these are the asymptotic
coefficients of $F_s$. No convergence of the Laurent series at a
finite real point is required.

Write $(x)_k=x(x+1)\cdots(x+k-1)$, $T_k(z)=(z-k/2)_k$, and
\[
 \mathcal R H(z)=\frac{H(z+1/2)+H(z-1/2)}{2z-1}.
\]
For positive even $n$ and $\ell=0,1,2$, let
\[
 H_{n,\ell}(z)=\frac{T_n(z)^5T_{2\ell}(z)}{(n!)^4},\qquad
 P_{n,\ell}=\mathcal R^nH_{n,\ell}.
\]
The exact divisibility in these operations, reflection
$P_{n,\ell}(1-z)=P_{n,\ell}(z)$, and degree $4n+2\ell$ are
the specialization of~\cite[Lemma 6.5]{KP}. Define a rational
functional by $\varphi_s(t^j)=c_{s,j}$ and put
\[
 Q_{n,s,\ell}(z)=\varphi_s\left(
                  \frac{P_{n,\ell}(z)-P_{n,\ell}(t)}{z-t}\right).
\]
This is the polynomial part of the formal product $P_{n,\ell}\Theta_s$.
The source's approximation theorem gives a remainder of order
$O(z^{-2n-2})$, with its constant depending on $n$.

Two other source inputs will be used explicitly. Its integer-basis
theorem~\cite[Theorem 8.1]{KP} writes $P_{n,\ell}$ as an integer
combination of $(z)_j(1-z)_j/(j!)^2$, $j\le2n+2$. Its rank
theorem~\cite[Theorem 7.1]{KP} says that the three columns
\[
       (P_{n,\ell},Q_{n,2,\ell},Q_{n,4,\ell})^T,
                       \qquad \ell=0,1,2,
\]
have a nonzero rational constant determinant. These are polynomial
statements; they do not require a $p$-adic evaluation.

\subsection{From a formal remainder to a real integral}
Put $\mu(y)=(2\cosh\pi y)^{-1}$ and $t_y=1/2+iy$. The Fourier
identity
\[
 \int_\R\mu(y)e^{iuy}\dd y=(2\cosh(u/2))^{-1}
\]
can be proved by a rectangular contour of height $i$; it is also the
Euler-moment identity used in Rivoal's real construction~\cite{Rivoal06}.
Fubini's theorem gives, for $\Re z>1/2$,
\[
 F_s(z)=(s-1)!\int_\R\frac{\mu(y)}{(z-t_y)^s}\dd y,
 \qquad
 \varphi_s(A)=\int_\R\mu(y)A^{(s-1)}(t_y)\dd y.
\]
For the first interchange an absolute majorant is
$u^{s-1}e^{-(\Re z-1/2)u}\mu(y)$. For the second identity compare
exponential moment generating functions and then coefficients.
Subtracting the polynomial part yields the exact remainder
\begin{align}
 E_{n,s,\ell}(z)&=P_{n,\ell}(z)F_s(z)-Q_{n,s,\ell}(z)\notag\\
 &=i^{s-1}\int_\R\mu^{(s-1)}(y)
                    \frac{P_{n,\ell}(t_y)}{z-t_y}\dd y.\label{eq:cauchy-remainder}
\end{align}
Integration by parts proves the last expression; exponential decay
of the weight kills every boundary term.

For an integer $N\ge1$, define $S_N(4)=\sum_{j<N}(-1)^j/(2j+1)^4$.
Splitting the defining beta series at $N$ gives
\[
 F_4(N+\tfrac12)=96(-1)^N\{\beta(4)-S_N(4)\}.
\]
Consequently $E_{n,4,\ell}(N+1/2)=b+a\beta(4)$, where
\begin{equation}\label{eq:pade-coefficients}
 a=96(-1)^NP_{n,\ell}(N+\tfrac12),\qquad
 b=-aS_N(4)-Q_{n,4,\ell}(N+\tfrac12).
\end{equation}
The rational tail in $b$ is part of the construction and must be cleared.
The rank theorem above ensures that at least one of these three forms
is nonzero at every such $n,N$: otherwise the last row of the constant
determinant would be $F_4(N+1/2)$ times the first.

For example, at $n=2,\ell=0,N=1$,
\[
 P(3/2)=31681/4096,\quad Q_4(3/2)=2223/256,
 \quad E(3/2)=-\frac{39}{256}\{4874\beta(4)-4817\}.
\]
The primitive integer form has absolute value about $3.12$. The small
prefactor in the displayed raw form provides no arithmetic saving.

\subsection{Transfer the difference operators to the kernel}
Estimating the large polynomial $P$ term by term loses its cancellations.
Moving its difference operators to the Cauchy kernel gives products with
a much clearer asymptotic balance. Define
\begin{align*}
 A_n(y)&=\prod_{j=1}^{n/2}\{y^2+(j-\tfrac12)^2\},\\
 B_\ell(y)&=\prod_{j=1}^{\ell}\{y^2+(j-\tfrac12)^2\},\\
 D_{n,N}(y)&=\prod_{j=0}^{n}\{(N-n/2+j)^2+y^2\}.
\end{align*}
An empty product is one.

\begin{proposition}\label{prop:product-integral}
For real $N\ge n/2$,
\[
 E_{n,4,\ell}(N+\tfrac12)
 =\frac{2(-1)^\ell}{(n!)^3}
   \int_0^\infty\frac{y\mu'''(y)A_n(y)^5B_\ell(y)}{D_{n,N}(y)}\dd y.
\]
At $N=n/2$ the integrand is interpreted continuously at $y=0$.
\end{proposition}
\begin{proof}
On the upward line $\Re t=1/2$, let $w_j(t)$ be the third derivative
of $\csc\pi(t+j/2)$. Oddness of $w_0$ about $1/2$ allows replacement
of the Cauchy kernel by its odd part. With
\[
 K_j(t,z)=\frac{j!(2t-1)}{(z-t-j/2)_{j+1}(z+t-1-j/2)_{j+1}},
\]
equation~\eqref{eq:cauchy-remainder} becomes
$E(z)=-(4i)^{-1}\int w_0(t)P(t)K_0(t,z)\dd t$.
Moving one difference operator from $P$ replaces $K_j$ by
\[
 \frac{K_j(t+1/2,z)}{2t}-\frac{K_j(t-1/2,z)}{2t-2}=K_{j+1}(t,z).
\]
For a direct check, put $a=z-t-(j+1)/2$ and
$b=z+t-(j+3)/2$. The numerator after taking a common denominator
is $j!\{(a+j+1)b-a(b+j+1)\}=(j+1)!(2t-1)$.

After $j$ transfers the remaining polynomial is divisible by $T_j^5$
by the source divisibility lemma. In the next transfer the vertical
lines shift to real parts zero and one. The poles of $w_{j+1}$ in
the crossed strips have order at most four and are canceled by zeros
of $T_{j+1}^5$. The apparent factors $2t$ and $2t-2$ have canceled
in the kernels. For $N>n/2$ the remaining kernel poles lie outside
the respective crossed strips. Exponential decay kills the horizontal
closing segments, so no residue is lost.

After $n$ transfers, put $t=1/2+iy$. Then
\[
 w_n=(-1)^{n/2}2i\mu'''(y),\quad
 T_n=(-1)^{n/2}A_n(y),\quad T_{2\ell}=(-1)^\ell B_\ell(y),
 \quad K_n=2iy n!/D_{n,N}(y).
\]
Substitution and pairing $y$ with $-y$ give the claimed normalization
and sign. As $N\downarrow n/2$, $\mu'''(y)=O(y)$ cancels the only
new zero denominator: $y^2/((N-n/2)^2+y^2)\le1$. The other factors
supply an integrable majorant near zero and infinity. Dominated
convergence gives the endpoint.
\end{proof}

\subsection{The actual exponential rate}
For $r\ge1/2$ define
\[
 \Phi_r(u)=3-\pi u+5\int_0^{1/2}\log(u^2+x^2)\dd x
                -\int_{r-1/2}^{r+1/2}\log(u^2+x^2)\dd x.
\]
The constant $3$ comes from $(n!)^{-3}$, the term $-\pi u$ from
the hyperbolic weight, and the two integrals from the numerator and
denominator products after setting $y=nu$. All powers of $n$ cancel.

\begin{theorem}\label{thm:pade-rate}
The function $\Phi_r$ has a unique maximum at $u_r>1/2$. Put
$\Lambda(r)=\Phi_r(u_r)$. If $n$ is even, $N\ge n/2$ is integral,
and $N/n\to r\ge1/2$, then for each $\ell=0,1,2$,
\[
 \lim\frac1n\log|E_{n,4,\ell}(N+\tfrac12)|=\Lambda(r).
\]
Its sign is eventually $(-1)^{\ell+1}$.
\end{theorem}
\begin{proof}
For $u>0$,
\begin{align*}
 \Phi_r'(u)&=-\pi+10\arctan(1/(2u))
             -2\arctan((r+1/2)/u)+2\arctan((r-1/2)/u),\\
 \Phi_r''(u)&=-\frac5{u^2+1/4}
       +\frac{2(u^2-ab)}{(u^2+a^2)(u^2+b^2)}
       \le-\frac3{u^2+1/4}<0,
\end{align*}
where $a=r-1/2$, $b=r+1/2$. The derivative is positive at $1/2$
and tends to $-\pi$ at infinity. This proves the unique maximum.

The sign issue in the integral is localized:
\[
 \mu'''(y)=\frac{\pi^3}{2}\sech(\pi y)\tanh(\pi y)
                         \{6\sech^2(\pi y)-1\}.
\]
It is positive up to the fixed value
$y_0=\operatorname{arcosh}(\sqrt6)/\pi$ and negative afterwards.
All the product factors are positive. The saddle region $y\sim nu_r$
therefore has one sign and is separated from the bounded region where
the sign is opposite.

Put $r_n=N/n$. On compact positive $u$ intervals, product Riemann
sums and Stirling's formula give $n\Phi_{r_n}(u)+O(\log n)$, uniformly
for $r_n$ in a fixed compact subset of $[1/2,\infty)$.
The product estimates retain this logarithmic error down to
$u=y_0/n$ on bounded intervals: bracket each
monotone Riemann sum by its endpoint sums, whose difference is
$O(\log n)$. The integral formula for $\Phi_r(u)$ is jointly continuous
for $r\ge1/2$ and $u\ge0$. Hence $r_n\to r$ gives
$\Phi_{r_n}\to\Phi_r$ uniformly on bounded $u$ intervals. After
division by $n$, the limiting profile is therefore $\Phi_r$;
no rate of convergence of $r_n$ is needed.

To bound $0\le y\le y_0$ separately, use
$A_n(y)\le C A_n(0)$, which follows from $\log(1+x)\le x$ and
the convergent sum of $(j-1/2)^{-2}$. Put $k=N-n/2$. For $k\ge1$
this part is at most
\[
 \frac{C A_n(0)^5}{(n!)^3(k)_{n+1}^2},
\]
and for $k=0$, after canceling $y^2$, at most $C A_n(0)^5/(n!)^5$.
The gamma identities for these products give upper rate $\Phi_r(0)$
uniformly even when $k/n\to0$. This is strictly below $\Phi_r(u_r)$.

For the infinite tail, use
$A_n(nu)\le n^n(u^2+1/4)^{n/2}$,
$D_{n,N}(nu)\ge(nu)^{2n+2}$ and $n!\ge(n/e)^n$.
Apart from polynomial factors, the tail is bounded by
\[
 \exp\{n[3-\pi u+\tfrac52\log(u^2+1/4)-2\log u]\}.
\]
For sufficiently large $u$ the bracket is at most $-\pi u/2$.
Laplace's estimate near the strict maximum now supplies the claimed
rate. The opposite-sign part has smaller rate, so cannot cancel it.
\end{proof}

The saddle and rate admit exact rational certification. If $w_r=u_r^2$,
then $w_r$ is the unique root greater than $1/4$ of
\[
 64w^3+(64r^2-16)w^2-(160r^2+20)w+20r^2-3=0.
\]
This follows by taking the real part of
$(u+i/2)^5(u+i(r-1/2))(u-i(r+1/2))$ in the saddle equation.
Descartes' rule and the signs at zero and $1/4$ select the larger
positive root. The arctangent terms cancel at the saddle, giving
\begin{align*}
 \Lambda(r)={}&\tfrac52\log(w_r+1/4)
 -(r+1/2)\log(w_r+(r+1/2)^2)\\
 &+(r-1/2)\log(w_r+(r-1/2)^2).
\end{align*}
For example $w_2=7/4$ and $\Lambda(2)=-2\log2$ exactly.

\subsection{Intersect two denominator bounds}
Let $d_j=\lcm(1,\ldots,j)$ and let a superscript ``odd'' mean the
odd part. Put
\[
 J=2n+2,\quad R_n^{\rm ar}=\gcd(n!,d_J)^{\rm odd},\quad
 L_{n,N}=\lcm(R_n^{\rm ar},d_{2N-1}^{\rm odd}).
\]
\begin{proposition}\label{prop:pade-clearing}
Both coefficients in~\eqref{eq:pade-coefficients} become integral after
multiplication by
\[
 M_{n,\ell,N}=2^{4n+2\ell-4+4v_2(n!)}L_{n,N}^4.
\]
For $N/n\to r\ge1/2$, its logarithmic rate is $8\log2+8r$.
\end{proposition}
\begin{proof}
First $(n!)^4P_{n,\ell}\in\Z[z]$. Expand the monic integer seed
$T_n^5T_{2\ell}$ in the monic basis $T_{n+2j}$ by successive leading
terms. Reflection preserves the even degrees and divisibility by $T_n$
forces the final remainder to vanish. The identity
$\mathcal R T_k=T_{k-1}$ proves the assertion. Recurrence
\eqref{eq:pade-moments} gives
$c_{4,0}=c_{4,1}=c_{4,2}=0$ and $2^{j-3}c_{4,j}\in\Z$ for $j\ge3$.
Thus the odd denominator of $Q_{n,4,\ell}(N+1/2)$ divides the odd
part of $(n!)^4$.

The factorial-basis calculation of~\cite[Sections 5 and 9]{KP}
gives a second bound. To state the needed arithmetic explicitly, its
coefficients satisfy, for $j\ge1$,
\[
 2b_j=\sum_{u+v=j-1}\frac{u!^2v!^2}{(u+1)(2u+1)!(2v+1)!},
 \qquad
 \Theta_4(z)=\sum_{j\ge1}\frac{b_j(2j+1)!}{(z)_{j+1}(1-z)_{j+1}}.
\]
One has $2d_j^2(2j+1)!b_j/(j!)^2\in\Z$: a summand multiplied by
$(2j+1)!/(j!)^2$ is
$2\binom{2j+1}{2u+2}(u!v!/j!)^2$, and $j!/(u!v!)$ divides $d_j$.
At each prime power its valuation contribution is zero or one because
$u+v=j-1$. Likewise $(j+1)\binom{k}{j+1}$ divides $d_k$.
Canceling the factorials against an integer basis term of $P$ shows
that the odd denominator of $Q(N+1/2)$ divides $(d_J^{\rm odd})^4$.
Half-integer factorial evaluations introduce no other odd primes;
indeed $2^{2k}(x)_k/k!\in\Z$ for half-integral $x$.
Intersecting the two bounds gives $(R_n^{\rm ar})^4$.

The polynomial $P$ has no odd denominator at a half-integer, by its
integer factorial basis. The finite tail $S_N(4)$ therefore adds
only $(d_{2N-1}^{\rm odd})^4$. Taking their lcm gives $L_{n,N}^4$.

For powers of two let $D=4n+2\ell$. A term $P_dc_{4,j}z^{d-j-1}$
of the polynomial part has valuation at least
$-4v_2(n!)-(j-3)-(d-j-1)\ge-4v_2(n!)-D+4$.
The coefficient of $\beta(4)$ includes $96$, with valuation five;
$S_N(4)$ has odd denominator. The stated power of two clears both.

Finally $d_n^{\rm odd}\mid R_n^{\rm ar}$, and any additional prime
power over $d_n$ has base at most $\sqrt J$. Its total logarithmic
cost is $O(\sqrt n\log n)$. Hence
$\log L_{n,N}=\log d_{\max(n,2N-1)}^{\rm odd}+O(\sqrt n\log n)$.
The prime number theorem and $v_2(n!)=n+O(\log n)$ prove the rate.
\end{proof}

\subsection{The gap that remains}
Combine the two proved rates and write
$\mathcal G(r)=8\log2+8r+\Lambda(r)$.
This is the logarithmic rate of the cleared forms; a negative value
would make them tend to zero.
\begin{center}
\begin{tabular}{@{}rrrr@{}}
\toprule
$r$ & Remainder rate & Clearing cost & Sum $\mathcal G(r)$\\
\midrule
$1/2$ & $-0.21025504$ & $9.54517744$ & $9.33492240$\\
$1$ & $-0.61293790$ & $13.54517744$ & $12.93223954$\\
$2$ & $-1.38629436$ & $21.54517744$ & $20.15888308$\\
\bottomrule
\end{tabular}
\end{center}
This is an asymptotic comparison for the explicit multiplier, not a fit
to finite values. Differentiation at the unique maximum gives
\[
 \Lambda'(r)=\log\frac{u_r^2+(r-1/2)^2}{u_r^2+(r+1/2)^2}.
\]
Its negative is $\int_{r-1/2}^{r+1/2}2x/(u_r^2+x^2)\dd x
\le1/u_r<2$. Therefore $\mathcal G'(r)>6$: moving the evaluation
point farther right makes this balance worse. The minimum is
$\mathcal G(1/2)=9.33492239999\ldots$.

This does not lower-bound the primitive forms. If the two cleared
coefficients have gcd $g_n$, their primitive value is $|M E|/g_n$.
Decay would require $\liminf n^{-1}\log g_n\ge\mathcal G(r)$.
No upper bound excluding that saving has been proved. Evaluation below
$N=n/2$, combinations of the three forms and modified seeds are also
outside this rate theorem. The finite tests and their precise ranges
are recorded in~\rec{R26}--\rec{R27}.
